% ============================================================================
% KAPITEL 9: Wirbelstruktur, Maxwell-Gravitation und Ω-Theorie
% ============================================================================

\chapter{Wirbelstruktur, Maxwell-Gravitation und \( \Omega \)-Theorie}
\label{ch:maxwell_gravitation}

\section{Einleitung}
\label{sec:mg_einleitung}

Die Weber-Gravitation (WG) enthält in ihrem geschwindigkeitsabhängigen Term eine implizite Wirbelstruktur, die bei geeigneter Interpretation direkt zu den Gleichungen der Maxwell-Gravitation (MG) führt. Die Dynamische Schwere-Trägheits-Theorie (DSTT) liefert die energetische Interpretation dieser Wirbel und vermittelt zwischen der differentiellen Kraftbeschreibung der WG und der integralen Feldbeschreibung der MG.

Die \( \Omega \)-Theorie ist die natürliche Feldtheorie, die aus dieser Maxwell-Gravitation emergiert. Sie basiert auf der Idee, dass \textbf{Rotation elementar ist, nicht Krümmung} – eine Alternative zur Raumzeitkrümmung der Allgemeinen Relativitätstheorie (ART).

\section{Grundlegende Axiome einer Wellenbeschreibung der Gravitation}
\label{sec:wellen_axiome}

Jeder physikalischen Wellenbeschreibung liegen fünf Axiome zugrunde, die auch der WDBT+ implizit zugrunde liegen:

\begin{enumerate}
    \item \textbf{Welle:} Jede Welle besteht aus einem Träger (Medium oder Feld) und einer Information. Die Information kann den Wert Null annehmen (reine Trägerwelle).
    \item \textbf{Information:} Information ist eine Störung der Welle – jede zeitliche oder räumliche Abweichung von der reinen, ungestörten Wellenform.
    \item \textbf{Kraft:} Kraft ist der negative Gradient der Energie über den Ort: \(\vec{F} = -\nabla E_{\text{pot}}\).
    \item \textbf{Informationsübertragung:} Informationsübertragung erfordert eine Störung, die sich mit endlicher, durch das Medium bestimmter Geschwindigkeit ausbreitet.
    \item \textbf{Form des Trägers:} Die periodische Funktion, die den Träger beschreibt, kann jede beliebige Gestalt annehmen; die Natur realisiert die energetisch optimale Form.
\end{enumerate}

\section{Die stehende Welle als Ursache der instantanen Kraft}
\label{sec:stehende_welle}

Zwei Massen \(M_1\) und \(M_2\) im Abstand \(d\) reflektieren die Wellen des Gravitationsfeldes. Zwischen ihnen entsteht durch Überlagerung von Hin- und Rückwelle eine stehende Welle:

\[
\Phi(z, t) = \Phi_1 \cos(kz - \omega t) + \Phi_2 \cos(kz + \omega t)
\]

Für gleiche Amplituden \(\Phi_1 = \Phi_2 = \Phi_0\) folgt:

\[
\Phi(z, t) = 2\Phi_0 \cos(kz) \cos(\omega t)
\]

Dies ist eine reine Trägerwelle ohne Informationsgehalt. Die Kraft auf eine Probemasse \(m\) ergibt sich aus dem Gradienten:

\[
F_z(z, t) = -m \frac{\partial \Phi}{\partial z} = 2m \Phi_0 k \sin(kz) \cos(\omega t)
\]

Die Kraft existiert \emph{instantan} an jedem Ort – die stehende Welle muss nicht erst „ankommen“. Im langwelligen Grenzfall (\(\lambda \gg d\), also \(k \to 0\)) geht diese Beschreibung in das Newton-Potential über:

\[
\Phi_{\text{eff}}(r) = -\frac{GM}{r}, \qquad \vec{F} = -\frac{GMm}{r^2} \hat{r}
\]

\section{Gravitationswellen als Störungen der stehenden Welle}
\label{sec:gravitationswellen_stoerung}

Wenn sich die Massen relativ zueinander bewegen (z. B. bei einer Verschmelzung), ändert sich die stehende Welle. Diese Änderung ist eine Störung des vorher stationären Zustands:

\[
\delta \Phi(\vec{r}, t) = \Phi(\vec{r}, t) - \Phi_{\text{stat}}(\vec{r})
\]

Diese Störung breitet sich als laufende Welle mit der endlichen Geschwindigkeit \(c\) aus:

\[
\square \delta \Phi = 0, \qquad \square = \frac{1}{c^2} \frac{\partial^2}{\partial t^2} - \nabla^2
\]

Damit sind die beobachteten Gravitationswellen (z. B. GW170817, siehe \cite{LIGO2017}) vollständig erklärbar – ohne dass die \emph{Kraft selbst} retardiert sein müsste.

\section{Der fundamentale Unterschied zur ART}
\label{sec:unterschied_art}

Die ART behauptet, dass sich die \emph{Kraft selbst} mit Lichtgeschwindigkeit ausbreitet. Dies ist aus Sicht der WDBT+ ein Kategorienfehler:

\begin{itemize}
    \item Die \textbf{Kraft} ist der Gradient einer stehenden Welle und wirkt instantan.
    \item Nur die \textbf{Störung} dieser Welle (also die Information „die Quelle bewegt sich“) breitet sich mit \(c\) aus.
\end{itemize}

Die beobachteten Gravitationswellen beweisen daher nicht die ART, sondern lediglich die endliche Ausbreitungsgeschwindigkeit von \emph{Veränderungen}. Die ART baut zudem einen kompensierenden Vektorpotential-Term ein, um die unphysikalische Konsequenz einer rein retardierten Gravitation (die das Sonnensystem destabilisieren würde) zu verhindern – ein mathematischer Trick, den die WDBT+ nicht benötigt.

\subsection{Konsequenzen für die Dynamik des Sonnensystems}
\label{sub:sonnensystem_dynamik}

Würde sich die Gravitation \emph{ohne} diesen Kompensationsterm mit \(c\) ausbreiten, so würde die Erde nicht der aktuellen Position der Sonne folgen, sondern ihrer Position vor etwa 8 Minuten. Die resultierende tangentiale Kraftkomponente wäre zwar klein (\(\sim 3 \times 10^{-8}\) der Radialkraft), aber sie würde über astronomische Zeiträume zu einer messbaren Drift führen. Das Sonnensystem wäre nicht stabil.

Da das Sonnensystem stabil ist, folgt daraus zwingend: Die Gravitation muss \emph{effektiv instantan} wirken. Die WDBT+ liefert mit der stehenden Welle die physikalisch konsistente Erklärung dafür, ohne die ART-Retardierungs-Tricks.

\section{Die Wirbelstruktur der Weber-Gravitation}
\label{sec:wirbel_wg}

\subsection{Der geschwindigkeitsabhängige Term}
\label{sub:tan_term_wg}

Der für die Wirbelstruktur entscheidende Term in der WG ist:

\[
\vec{F}_{\text{tan}} = -\frac{GMm}{r^2} \cdot \frac{2(\dot{r} \cdot \vec{v})}{c^2} \vec{v}
\]

Dieser Term hat folgende Eigenschaften:

\begin{itemize}
    \item Er ist nicht-zentral: seine Richtung hängt von \( \vec{v} \) ab, nicht nur von \( \vec{r} \)
    \item Er ist geschwindigkeitsabhängig (linear in \( \vec{v} \) für jede Komponente)
    \item Er besitzt eine nichtverschwindende Rotation (Wirbelstärke)
    \item Er ist analog zum geschwindigkeitsabhängigen Term in der Weber-Elektrodynamik
\end{itemize}

\subsection{Die Wirbeldichte der WG}
\label{sub:wirbeldichte_wg}

Definiere die Wirbeldichte (Rotation) des Kraftfeldes:

\[
\vec{\omega} = \nabla \times \vec{F}_{\text{WG}} = \nabla \times (\vec{F}_{\text{rad}} + \vec{F}_{\text{tan}})
\]

Der radiale Anteil \( \vec{F}_{\text{rad}} \) ist ein Zentralfeld und hat keine Rotation (\( \nabla \times \vec{F}_{\text{rad}} = 0 \)). Die gesamte Wirbelstärke kommt vom tangentialen Anteil:

\[
\vec{\omega} = \nabla \times \vec{F}_{\text{tan}} = -\frac{2GMm}{c^2} \nabla \times \left( \frac{(\dot{r} \cdot \vec{v}) \vec{v}}{r^2} \right)
\]

Nach Ausführung der Rotation ergibt sich:

\[
\boxed{\vec{\omega} = \frac{6GMm}{c^2 r^3} (\hat{r} \cdot \vec{v})(\hat{r} \times \vec{v})}
\]

Dies ist die fundamentale Wirbelstruktur der WG. Sie hat folgende Eigenschaften:

\begin{itemize}
    \item \( \vec{\omega} \) ist proportional zu \( 1/r^3 \) – wie ein Dipolfeld
    \item \( \vec{\omega} \) ist senkrecht zu \( \vec{r} \) und \( \vec{v} \) (wegen \( \hat{r} \times \vec{v} \))
    \item \( \vec{\omega} \) verschwindet, wenn \( \vec{v} \parallel \hat{r} \) (rein radiale Bewegung) oder \( \vec{v} \perp \hat{r} \) (reine Tangentialbewegung mit \( \hat{r} \cdot \vec{v} = 0 \))
    \item Maximale Wirbelstärke bei \( \hat{r} \cdot \vec{v} = \pm |\vec{v}|/\sqrt{2} \) (45°-Bahnneigung)
\end{itemize}

\section{Von der WG zur Maxwell-Gravitation}
\label{sec:wg_zu_mg}

\subsection{Identifikation des Vektorpotentials}
\label{sub:vektorpotential_mg}

Vergleiche den tangentialen Term \( \vec{F}_{\text{tan}} \) mit der Lorentz-Kraft \( \vec{F} = q(\vec{E} + \vec{v} \times \vec{B}) \). Der Term \( 2(\dot{r} \cdot \vec{v}) \vec{v} \) legt nahe, ein effektives Vektorpotential einzuführen:

\[
\vec{A}_g = \frac{2GM}{c^2 r} \vec{v}_q
\]

wobei \( \vec{v}_q \) die Geschwindigkeit der Quelle ist.

\subsection{Das gravito-magnetische Feld}
\label{sub:gravito_magnetisch}

Definiere das gravito-magnetische Feld als Rotation des Vektorpotentials:

\[
\vec{B}_g = \nabla \times \vec{A}_g
\]

Für eine Punktmasse mit konstanter Geschwindigkeit \( \vec{v}_q \) ergibt sich:

\[
\vec{B}_g = \frac{2G}{c^2} \cdot \frac{\vec{v}_q \times \vec{r}}{r^3}
\]

Dies ist exakt der gravito-magnetische Dipolterm, der aus der ART als Lense-Thirring-Effekt bekannt ist (siehe \cite{Lense1918}).

\subsection{Das gravito-elektrische Feld}
\label{sub:gravito_elektrisch}

Definiere das gravito-elektrische Feld aus dem radialen Anteil der WG:

\[
\vec{E}_g = -\frac{GM}{r^2} \left[ 1 - \frac{\dot{r}^2}{c^2} + \beta \frac{r\ddot{r}}{c^2} \right] \hat{r} + \text{Terme aus } \vec{A}_g
\]

Im statischen Limes (\( \dot{r} = 0, \ddot{r} = 0 \)) vereinfacht sich dies zu:

\[
\vec{E}_g = -\frac{GM}{r^2} \hat{r}
\]

\section{Die Maxwell-Gleichungen der Gravitation}
\label{sec:mg_gleichungen}

Mit diesen Definitionen ergeben sich die Maxwell-Gleichungen der Gravitation:

\[
\begin{aligned}
\nabla \cdot \vec{E}_g &= -4\pi G\rho + \mathcal{O}(v^2/c^2) \\
\nabla \cdot \vec{B}_g &= 0 \\
\nabla \times \vec{E}_g &= -\frac{\partial \vec{B}_g}{\partial t} + \mathcal{O}(v^2/c^2) \\
\nabla \times \vec{B}_g &= \frac{4\pi G}{c^2} \vec{j} + \frac{1}{c^2} \frac{\partial \vec{E}_g}{\partial t} + \mathcal{O}(v^2/c^2)
\end{aligned}
\]

Im Grenzfall verschwindender Korrekturen (\( v/c \to 0 \)) gehen sie in die bekannten Maxwell-Gleichungen der Gravitation über. Diese Gleichungen sind \textbf{glatt} – sie enthalten nur gewöhnliche Ableitungen (siehe auch Kapitel~\ref{ch:fraktale_maxwell} für die Erweiterung auf fraktale Quellen).

\section{Die Rolle der DSTT als Vermittler}
\label{sec:dstt_vermittler_mg}

Die DSTT liefert die energetische Interpretation der Wirbelstruktur. Die Zustandsvariable

\[
\beta_{\text{DSTT}}(t) = \frac{E_B}{E} = \frac{\frac{1}{2}mr^2\dot{\phi}^2}{\frac{1}{2}m\dot{r}^2 + \frac{1}{2}mr^2\dot{\phi}^2}
\]

beschreibt die Aufteilung der Energie in radiale und tangentiale Komponenten (siehe Kapitel~\ref{ch:dstt_weber_kraefte}). Diese Aufteilung ist direkt mit der Wirbeldichte verknüpft:

\[
|\vec{\omega}| = \frac{2GM}{c^2 r^2} \cdot \frac{\dot{\beta}_{\text{DSTT}}}{1 - \dot{\beta}_{\text{DSTT}}} \cdot \frac{r\dot{\phi}}{\dot{r}} + \mathcal{O}(v^2/c^2) \quad (\text{für } \dot{r} \neq 0)
\]

Diese Beziehung zeigt:

\begin{itemize}
    \item Die Wirbeldichte ist proportional zu \( \dot{\beta}_{\text{DSTT}} \) – die zeitliche Änderung der Energieaufteilung erzeugt Wirbel
    \item Für \( \dot{\beta}_{\text{DSTT}} = 0 \) (Elektrodynamik) verschwindet die Wirbeldichte im Rahmen dieser Näherung
    \item Die DSTT verbindet die \textbf{differentielle} Beschreibung (Wirbel) mit der \textbf{integralen} Beschreibung (Energieaufteilung)
\end{itemize}

\section{Die \( \Omega \)-Theorie der Gravitation}
\label{sec:omega_theorie}

Die \( \Omega \)-Theorie ist die natürliche Feldtheorie, die aus der Maxwell-Gravitation emergiert. Sie basiert auf der Idee, dass \textbf{Rotation elementar ist, nicht Krümmung}.

\subsection{Axiome der \( \Omega \)-Theorie}
\label{sub:omega_axiome}

\subsubsection{Axiom 1: Nicht-Lokalität und Instantaneität}
\label{sub:omega_axiom1}

Die Gravitationswirkung ist nicht-lokal und instantan. Es gibt keine Ausbreitung mit endlicher Geschwindigkeit; die Wechselwirkung ist simultan. Dies ist direkte Konsequenz der feldlosen Weber-Wechselwirkung der WDBT+ (siehe Abschnitte~\ref{sec:stehende_welle} und~\ref{sec:unterschied_art}).

\subsubsection{Axiom 2: Das Rotationsfeld}
\label{sub:omega_axiom2}

Die Gravitation wird durch ein fundamentales, axiales Vektorfeld \( \vec{\Omega}(\vec{r}, t) \) beschrieben. Dieses Feld kodiert die lokale Rotationsdichte des Raumes und ist die Ursache für alle geschwindigkeitsabhängigen (tangentialen) Trägheitseffekte.

\subsubsection{Axiom 3: Aufspaltung von Ursache und Wirkung}
\label{sub:omega_axiom3}

Die Schwere (Ursache) und die Trägheit (Wirkung) sind fundamental verschieden. Die radiale Anziehung wird durch ein skalares Potential \( \Phi(\vec{r}, t) \) beschrieben, die tangentialen Trägheitskräfte durch das Rotationsfeld \( \vec{\Omega}(\vec{r}, t) \).

\subsubsection{Axiom 4: DSTT-Kompatibilität}
\label{sub:omega_axiom4}

Die Theorie ist so konstruiert, dass aus ihr die Existenz einer Zustandsvariablen \( \beta(\vec{r}, t) \) folgt, die das lokale Verhältnis von tangentialer zu gesamter Trägheit angibt, und für die \( \dot{\beta} > 0 \) gilt. Dies verleiht der Gravitation einen intrinsischen Zeitpfeil.

\subsection{Die Feldgleichungen der \( \Omega \)-Theorie}
\label{sub:omega_feldgleichungen}

Die Felder werden nicht-lokal und instantan aus den Quellen erzeugt. Im Vakuum (außerhalb der Quellen) gehorchen sie formal den Maxwell-Gleichungen:

\[
\begin{aligned}
\nabla \cdot \vec{\Omega} &= 0 \\
\nabla \times \vec{\Omega} &= \frac{4\pi G}{c^2} \vec{j} + \frac{1}{c^2} \frac{\partial}{\partial t} (-\nabla \Phi) \\
\nabla \cdot (-\nabla \Phi) &= -4\pi G\rho \\
\nabla \times (-\nabla \Phi) &= -\frac{\partial \vec{\Omega}}{\partial t}
\end{aligned}
\]

Es ist jedoch entscheidend zu betonen, dass diese Differentialgleichungen hier nicht fundamental sind, sondern eine alternative Schreibweise der nicht-lokalen Integralgleichungen darstellen. Die wahre Natur der Theorie ist die der instantanen, integralen Wechselwirkung.

\subsection{Omega-Wellen: Gravitationswellen in der \( \Omega \)-Theorie}
\label{sub:omega_wellen}

Die \( \Omega \)-Theorie sagt Gravitationswellen voraus – aber sie unterscheiden sich fundamental von denen der ART:

\begin{itemize}
    \item \textbf{Fundamentale Natur:} Omega-Wellen sind instantan korrelierte, kollektive Schwingungen des \( \vec{\Omega} \)-Feldes und der Massenverteilung.
    \item \textbf{Ausbreitung:} Keine Ausbreitungsverzögerung (fundamental), aber durch die fraktale Raumdimension \( D \) emergiert auf großen Skalen eine effektive Retardierung mit \( c \) (siehe Anhang~\ref{app:spektrale_dimension}).
    \item \textbf{Lokalität:} Nicht-lokal (fundamental), aber effektiv lokal auf großen Skalen.
    \item \textbf{Polarisation:} Axial (wirbelartig), nicht transversal wie bei ART-Wellen.
\end{itemize}

Die effektive Wellengleichung auf makroskopischen Skalen lautet:

\[
\square \vec{\Omega}_{\text{eff}} = \frac{4\pi G}{c^2} \vec{j}_{\text{eff}} + \mathcal{O}\left(\frac{1}{D-3}\right)
\]

Die ART ist in dieser Sicht der Grenzfall \( D \to 3 \).

\subsection{Testbare Vorhersagen}
\label{sub:omega_vorhersagen}

Die Theorie macht spezifische, von der ART abweichende Vorhersagen (siehe auch Kapitel~\ref{ch:testbare_vorhersagen}):

\begin{enumerate}
    \item \textbf{Frequenzabhängige Geschwindigkeit:}
    \[
    v_{\text{phase}}(\omega) = c \left( 1 + \alpha \left( \frac{\omega}{\omega_0} \right)^{D-3} + \dots \right)
    \]
    Mit \( D \approx 2,704 \) ist \( D-3 \approx -0,296 \), also eine schwache Dispersion.

    \item \textbf{Fraktales Spektrum:} Die Leistungsverteilung skaliert mit \( P(\omega) \sim \omega^{-\beta} \), \( \beta \approx 0,7 \).

    \item \textbf{Nicht-lokale Korrelationen:} Bei Verschmelzungsereignissen sollten instantane Korrelationen zwischen weit entfernten Detektoren auftreten.

    \item \textbf{Abweichungen bei hohen Frequenzen:} Oberhalb einer kritischen Frequenz weichen die Signale systematisch von ART-Vorhersagen ab.
\end{enumerate}

\section{Krümmung der ART wird zur Rotation in der \( \Omega \)-Theorie}
\label{sec:kruemmung_zu_rotation}

Die ART beschreibt Gravitation als Krümmung der Raumzeit. Die \( \Omega \)-Theorie beschreibt Gravitation als Rotationsfeld \( \vec{\Omega} \). Der Übergang wird durch die DSTT und die Weber-Kräfte vermittelt:

\begin{itemize}
    \item Die \textbf{Krümmung} der ART entspricht in der \( \Omega \)-Theorie der \textbf{Wirbelstärke} des Rotationsfeldes.
    \item Die \textbf{Geodätengleichung} der ART wird durch die Bewegungsgleichung \( \vec{F} = m(-\nabla \Phi + 2\vec{v} \times \vec{\Omega} - \partial\vec{A}/\partial t) \) ersetzt.
    \item Die \textbf{Raumzeit-Metrik} der ART ist emergent – sie entsteht aus der fraktalen Informationsstruktur der IWT (Teil I, siehe Kapitel~\ref{ch:emergenz_raum_zeit} und Anhang~\ref{app:kontinuumslimes}).
\end{itemize}

Damit ist die \( \Omega \)-Theorie eine vollwertige Alternative zur ART: Sie vermeidet Singularitäten, kommt ohne dunkle Materie und dunkle Energie aus, und sie bewahrt die Nicht-Lokalität und Irreversibilität der fundamentalen WDBT+.

\section{Zusammenfassung der Zusammenhänge}
\label{sec:mg_zusammenhaenge}

Die folgende Tabelle fasst die verschiedenen Ebenen und ihre Verbindungen zusammen:

\begin{tabular}{|l|l|l|}
\hline
\textbf{Ebene} & \textbf{Größe} & \textbf{Zusammenhang} \\
\hline
WG (Kraft) & \(F_{\text{tan}} \propto (\dot{r} \cdot \vec{v})\vec{v}\) & Enthält Wirbelterm \\
Wirbeldichte & \(\vec{\omega} = \nabla \times \vec{F}_{\text{WG}}\) & Charakterisiert Rotation \\
Vektorpotential & \(\vec{A}_g = \frac{2GM}{c^2 r} \vec{v}_q\) & Emergiert aus WG \\
gravito-magnetisches Feld & \(\vec{B}_g = \nabla \times \vec{A}_g\) & Lense-Thirring-Effekt \\
Maxwell-Gleichungen & \(\nabla \cdot \vec{B}_g = 0,\; \nabla \times \vec{B}_g = \frac{4\pi G}{c^2} \vec{j} + \ldots\) & Emergieren im Kontinuumslimes \\
DSTT & \(\beta(t) = E_B/E\) & Energetische Interpretation der Wirbel \\
\(\dot{\beta}\)-Relation & \(|\vec{\omega}| \propto \dot{\beta}/(1-\beta)\) & Verbindet Wirbel und Energieaufteilung \\
\hline
\end{tabular}

\section{Fazit}
\label{sec:mg_fazit}

Die Weber-Gravitation ist mehr als eine bloße Kraftformel. Sie enthält in ihrer Struktur bereits die Keimzelle einer vollständigen Feldtheorie der Gravitation:

\begin{itemize}
    \item Die Gravitation wird als stehende Welle zwischen Massen beschrieben. Die Kraft ist der Gradient dieser Welle und wirkt instantan. Veränderungen (Gravitationswellen) breiten sich mit \( c \) aus.
    \item Newtons Gravitationsgesetz emergiert als statischer Grenzfall. Die Periheldrehung des Merkur ist bereits in instantanen Theorien (Weber, \cite{Tisserand1894}) enthalten – kein Privileg der ART.
    \item Der geschwindigkeitsabhängige Term der WG erzeugt eine Wirbeldichte, die bei geeigneter Interpretation direkt zum gravito-magnetischen Feld der Maxwell-Gravitation führt.
    \item Die DSTT erweist sich als der natürliche Vermittler zwischen diesen Beschreibungsebenen: Sie liefert die energetische Interpretation der Wirbel und verbindet die differentielle Kraftbeschreibung mit der integralen Feldbeschreibung.
    \item Die \( \Omega \)-Theorie formuliert Gravitation als Rotationsfeld \( \vec{\Omega} \), nicht als Raumzeitkrümmung. Die Krümmung der ART wird zur Wirbelstärke des \( \vec{\Omega} \)-Feldes.
\end{itemize}

WG, DSTT und MG sind keine unabhängigen Theorien, sondern verschiedene Aspekte ein und derselben physikalischen Realität – einer Realität, in der Gravitation nicht durch gekrümmte Raumzeit, sondern durch direkte Teilchenwechselwirkungen mit inhärenter Wirbelstruktur vermittelt wird.

Die ART ist eine hinreichende, aber keine notwendige Beschreibung. Ihre scheinbaren Erfolge beruhen auf einer Verwechslung von Kraft und Informationsübertragung. Die WDBT+ ist damit nicht nur konsistent mit allen bekannten Beobachtungen (Periheldrehung, Lichtablenkung, Gravitationswellen, Stabilität des Sonnensystems), sondern auch physikalisch anschaulicher und frei von den metaphysischen Annahmen einer gekrümmten Raumzeit.